\section{The mass closure}
\label{sec:closure}

\subsection{Decomposition}

\begin{definition}[Mass budget]
The gross mass is partitioned as
\begin{equation}
  m = \mfix + \mstr + \mprop + \mbat ,
  \label{eq:budget}
\end{equation}
where $\mfix$ is the payload and avionics that do not scale with aircraft size
(display, driver electronics, computer, camera, gimbal), $\mstr$ is the
structure, $\mprop$ is motors, speed controllers and propellers, and $\mbat$ is
the battery pack.
\end{definition}

\Cref{eq:budget} is an identity. It becomes an equation, and a nontrivial one,
once the three terms on the right that depend on $m$ are made explicit. The
essential point is that they depend on $m$ in different ways, and lumping them
together, as one is tempted to do by writing a single ``structural fraction'',
destroys the structure of the problem.

\subsection{Propulsion mass: sized by peak thrust}

\begin{definition}[Thrust margin]
The propulsion system is sized so that the maximum total thrust is
\begin{equation}
  T_{\max} = \lambda\, m g, \qquad \lambda > 1 .
  \label{eq:lambda}
\end{equation}
$\lambda$ is the ratio of maximum thrust to weight; $\lambda = 2$ means the
vehicle hovers at half throttle.
\end{definition}

\begin{model}[Specific thrust]
\label{mod:sm}
Let $S_m \coloneqq T_{\max}/\mprop$ be the specific thrust of the complete
propulsion system, in \si{\newton\per\kilo\gram}. Then
\begin{equation}
  \mprop = \frac{T_{\max}}{S_m} = \frac{\lambda g}{S_m}\, m
         \eqqcolon \gamma_m\, m .
  \label{eq:mprop}
\end{equation}
\end{model}

Propulsion mass is therefore \emph{linear} in gross mass. This is worth
emphasising because it is the coupling most people worry about first, and it is
benign: it consumes a fixed fraction of the mass budget and produces a runaway
only in the degenerate case $\gamma_m \ge 1$.

\Cref{mod:sm} is a lumped model. \Cref{sec:components} replaces it with motor
mass derived from the torque the rotor demands, at which point $S_m$ becomes an
output; the engine reports both so the assumption can be audited.

\subsection{Structural mass}

\begin{model}[Structural fraction]
\label{mod:fs}
\begin{equation}
  \mstr = f_s\, m, \qquad f_s \in (0,1) .
  \label{eq:mstr}
\end{equation}
\end{model}

This is the weakest model in the algebraic layer. It is not a law; it is an
observation that across a class of similar vehicles the structure tends to a
roughly constant fraction of gross mass, typically $f_s \in [0.15,0.25]$. It is
retained because it keeps the closure in closed form, which is what makes
\cref{sec:fold} possible. \Cref{sec:components} derives $\mstr$ from beam theory
instead and the two are compared in \cref{sec:results}.

\subsection{Battery mass: sized by mission energy}

\begin{model}[Battery]
\label{mod:battery}
Let $\eb$ be the specific energy of the pack in \si{\joule\per\kilo\gram} that
the mission may use: the cell value reduced by the depth of discharge and by the
landing reserve $\varsigma$, so $\eb = e_{\mathrm{cell}}\,\mathrm{DoD}\,(1-\varsigma)$.
Where the reserve must be visible, as in \cref{sec:codesign}, it is written out
and $\eb$ denotes the value before the reserve is removed. For a mission of
duration $t_f$ flown at total electrical power $P_{\mathrm{tot}}$,
\begin{equation}
  \mbat = \frac{E_{\mathrm{mission}}}{\eb}
        = \frac{1}{\eb}\int_0^{t_f} P\big(x(t),u(t)\big)\dd t
  \;\xrightarrow[\text{hover}]{}\;
  \frac{P_{\mathrm{tot}}\,t_f}{\eb} .
  \label{eq:mbat}
\end{equation}
\end{model}

The integral form is the honest one and is what \cref{sec:codesign} uses; the
right-hand form is its evaluation under the assumption that the whole mission is
spent hovering, and is what makes the algebra tractable. The ratio between them
is measured in \cref{sec:results} and is between $1.00$ and $1.10$ depending on
the mission, which retrospectively justifies the substitution for sizing
purposes but not for endurance claims.

\subsection{The master equation}

Substituting \cref{eq:mprop,eq:mstr,eq:mbat,eq:ptot} into \cref{eq:budget}:
\begin{equation}
  m = \mfix + f_s m + \gamma_m m
      + \frac{t_f}{\eb}\left[
        \frac{(mg)^{3/2}}{\FM\,\eta\,\sqrt{2\rho A}} + \Paux \right] .
  \label{eq:closure-raw}
\end{equation}
Collecting terms gives the central equation of the algebraic layer.

\begin{definition}[Closure coefficients]
\label{def:coeffs}
\begin{align}
  \alpha &\coloneqq 1 - f_s - \frac{\lambda g}{S_m}
          = 1 - f_s - \gamma_m
          &&\text{(thrust coupling, dimensionless)}, \label{eq:alpha}\\
  \beta  &\coloneqq \frac{t_f\, g^{3/2}}{\eb\,\FM\,\eta\,\sqrt{2\rho A}}
          &&\text{(energy coupling, \si{\kilo\gram^{-1/2}})}, \label{eq:beta}\\
  \Mzero &\coloneqq \mfix + \frac{\Paux\, t_f}{\eb}
          &&\text{(effective fixed load, \si{\kilo\gram})} . \label{eq:M0}
\end{align}
\end{definition}

\begin{theorem}[The closure equation]
\label{thm:closure}
Under \cref{mod:disk,mod:kappa,mod:sm,mod:fs,mod:battery} with fixed disk area,
the gross mass of a feasible design satisfies
\begin{equation}
  \boxed{\;\alpha\, m - \beta\, m^{3/2} = \Mzero\;}
  \label{eq:closure}
\end{equation}
and conversely any positive root of \cref{eq:closure} determines a mass budget
consistent with \cref{eq:budget}.
\end{theorem}

\begin{proof}
Rearranging \cref{eq:closure-raw},
\[
  m\left(1 - f_s - \gamma_m\right)
  - \frac{t_f g^{3/2}}{\eb \FM \eta \sqrt{2\rho A}}\, m^{3/2}
  = \mfix + \frac{\Paux t_f}{\eb},
\]
where the $m^{3/2}$ term arises because
$(mg)^{3/2} = g^{3/2} m^{3/2}$. Identifying the three groups with
\cref{eq:alpha,eq:beta,eq:M0} gives \cref{eq:closure}. The converse holds
because every step is an equivalence for $m>0$.
\end{proof}

\begin{remark}[Reading the coefficients]
$\alpha$ is the fraction of gross mass left over after structure and propulsion
have taken their cut; it is a pure number, must be positive for any design to
exist, and is easy to keep positive. $\beta$ carries the mission: it is
proportional to endurance and inversely proportional to specific energy,
figure of merit, electrical efficiency and the square root of disk area. It
multiplies the nonlinear term, and it is where the design dies. $\Mzero$ is
what must be carried: note that the auxiliary power appears here, not in
$\beta$, because it is independent of aircraft mass, so a display that consumes
power converts directly into fixed load through the battery needed to run it.
\end{remark}

\begin{example}
For a 13.3-inch panel with $\mfix = \SI{0.762}{\kilo\gram}$,
$f_s = 0.20$, $\lambda = 2$, $S_m = \SI{120}{\newton\per\kilo\gram}$,
$\FM = 0.65$, $\eta = 0.75$, $\eb = \SI{900}{\kilo\joule\per\kilo\gram}$,
$t_f = \SI{3600}{\second}$ and four \SI{0.40}{\metre} rotors
($A = \SI{0.503}{\metre\squared}$), \cref{def:coeffs} gives
$\alpha = 0.6366$, $\beta = \SI{0.2271}{\kilo\gram^{-1/2}}$,
$\Mzero = \SI{0.762}{\kilo\gram}$. \Cref{sec:fold} shows that this
configuration has no solution.
\end{example}

\subsection{Generalised disk-area scaling}

\Cref{cor:hover} assumed the disk area fixed while mass varies. That is an
assumption about the design family, not a law. Suppose instead the rotors are
allowed to grow with the aircraft,
\begin{equation}
  A(m) = A_0 \left(\frac{m}{m_0}\right)^{p} ,
  \label{eq:areascaling}
\end{equation}
anchored so that $A(m_0) = A_0$ reproduces a stated rotor at a stated mass.

\begin{proposition}[Generalised power law]
\label{prop:plaw}
Under \cref{eq:areascaling} the hover propulsive power obeys
\begin{equation}
  \Pelec \propto m^{q}, \qquad q = \frac{3-p}{2} ,
  \label{eq:qexp}
\end{equation}
and the closure becomes
\begin{equation}
  \alpha m - \beta_q m^{q} = \Mzero,
  \qquad
  \beta_q = \frac{t_f\,g^{3/2}\,m_0^{p/2}}
                 {\eb\,\FM\,\eta\,\sqrt{2\rho A_0}} .
  \label{eq:closure-q}
\end{equation}
\end{proposition}

\begin{proof}
Substituting \cref{eq:areascaling} into \cref{eq:pelec} with $T=mg$,
\[
  \Pelec = \frac{(mg)^{3/2}}{\FM\eta\sqrt{2\rho A_0}}
           \left(\frac{m}{m_0}\right)^{-p/2}
         = \frac{g^{3/2} m_0^{p/2}}{\FM\eta\sqrt{2\rho A_0}}\; m^{(3-p)/2}.
\]
Multiplying by $t_f/\eb$ and inserting into \cref{eq:budget} as before gives
\cref{eq:closure-q}.
\end{proof}

Three cases matter and are treated in \cref{sec:fold,sec:scaling}:
\begin{itemize}[nosep]
  \item $p=0$: rotors fixed, $q=3/2$, the case of \cref{thm:closure};
  \item $p=1$: constant disk loading, $q=1$, the closure becomes linear;
  \item $p>1$: rotors grow faster than mass, $q<1$, the nonlinearity is
        sublinear and inverts its role entirely.
\end{itemize}

\begin{corollary}[Constant disk loading]
\label{cor:constdl}
If the rotors are sized to hold $\mathrm{DL} = mg/A$ constant, then by
\cref{eq:dl} the propulsive power is exactly linear in mass,
$\Pelec = c_P m$ with
\begin{equation}
  c_P = \frac{g}{\FM\,\eta}\sqrt{\frac{\mathrm{DL}}{2\rho}} ,
  \label{eq:cP}
\end{equation}
the battery mass fraction is $\gamma_b = c_P t_f/\eb$, and the closure
collapses to
\begin{equation}
  m = \frac{\Mzero}{1 - f_s - \gamma_m - \gamma_b} ,
  \label{eq:constdl}
\end{equation}
which has a positive solution if and only if
\begin{equation}
  f_s + \gamma_m + \gamma_b < 1 .
  \label{eq:constdlfeas}
\end{equation}
\end{corollary}

\begin{proof}
Immediate from \cref{eq:dl} with $T=mg$ and \cref{eq:closure-q} at $q=1$,
$\beta_1 = \gamma_b$.
\end{proof}

The contrast between \cref{eq:closure} and \cref{eq:constdl} is the central
structural fact of the sizing problem, and it is the subject of the next
section. In the first, feasibility fails discontinuously at a fold. In the
second, mass diverges continuously as the fractions approach unity. They are
different kinds of impossibility.
